\documentclass[11pt]{article}
\usepackage[a4paper,margin=1in]{geometry}
\usepackage[T1]{fontenc}
\usepackage{amsmath,amssymb,booktabs,tabularx}
\usepackage[authoryear,round]{natbib}
\usepackage[hidelinks]{hyperref}
\usepackage{siunitx}
\sisetup{group-digits=integer,group-separator={,},group-minimum-digits=4}
\newcommand{\conc}{\mathrm{C}}
\newcommand{\spr}{\mathrm{S}}
\title{Patient Shortage in Histopathology:\\Does Unequal Patient Influence Explain the Damage?}
\author{Yannik Queisler}
\date{September 10, 2026}

\begin{document}
\raggedbottom
\widowpenalty=10000
\clubpenalty=10000
\setlength{\emergencystretch}{2em}
\maketitle
\begin{abstract}
\noindent Reducing the number of training patients can impair classification even when class patch counts remain fixed. The preceding experiments found that several learning methods and alternative classifiers recover little of this damage. This prospective experiment asks whether unequal patient influence contributes to the remaining deficit. Multinomial logistic regression is compared under ordinary patch-average and patient-average training, using identical concentrated and spread patient allocations from BRACS and TCGA-UT. Patient averaging gives each contributing patient equal loss mass within a class while preserving that class's total influence. Frozen features, existing splits, and selected regularization remain fixed. The comparison separates the gain from equalizing patient influence from the benefit of broader coverage that remains afterwards. A residual coverage benefit would identify a limitation of this correction, without establishing which missing variation causes the loss.
\end{abstract}
\tableofcontents
\clearpage

\section{Motivation}
\label{sec:motivation}
The signal taxonomy of \citet{queisler2026imbalanceMethods} distinguishes nominal support, the number of labelled examples, from independent support, the number of patients supplying them. Patches from one patient share biological and acquisition characteristics, so their count need not reflect the amount of independent evidence. Within-class diversity describes the variation observed under a specified representation. Together, these concepts explain why a large patch collection can still provide limited exposure to differences between patients.

The controlled shortage benchmark of \citet{queisler2026benchmarkResults} reduced patient support while holding class patch counts fixed. Discrimination losses reached 2.62--3.03 percentage points on BRACS and 4.76--4.83 points on TCGA-UT, and tested methods left these losses largely unrecovered. The classifier comparison of \citet{queisler2026patientShortage} evaluated the original multilayer perceptron alongside logistic regression and cosine nearest neighbours. Under concentrated support, logistic regression improved patient-macro balanced accuracy by only 0.19 points on BRACS and 0.48 points on TCGA-UT. Broader coverage benefited logistic regression by 2.44 points (95\% interval: 0.31 to 4.98) and 4.70 points (4.04 to 5.34), respectively. These findings support a coverage limitation within the tested procedures and motivate examining how patient shortage affects learning.

\emph{Unequal patient influence} arises because ordinary training assigns equal weight to patches. Within a class, a patient supplying more patches consequently contributes more of the training loss. When few patients supply the available patches, the fitted classifier may depend strongly on the characteristics of those with the largest contributions. Those characteristics may be less representative of patients encountered at evaluation. Patient-average training redistributes the loss so that each contributing patient has equal weight within a class. A gain on the same training patches would support unequal influence as a contributor to the observed damage. Here, influence refers to loss weight; the observations themselves still determine how each patient affects the fitted decision rule.

Broader patient coverage also exposes the classifier to additional variation within each class, including differences in tissue appearance and acquisition conditions. Even when all observed patients receive equal weight, a small patient set may represent this variation incompletely. A \emph{residual coverage limitation} is the performance benefit of broader coverage that remains after patient averaging. This residual describes how much of the coverage deficit survives the correction. Unequal influence and incomplete coverage can therefore contribute together: weighting may recover part of the loss while additional patients remain beneficial.

The central question is whether equalizing patient influence reduces the established coverage deficit. The comparison measures both the gain under concentrated support and the remaining benefit of spread support. The benchmark's round-robin sampling and contribution caps already limit unequal influence \citep{queisler2026benchmarkProtocol}. The observed contribution distribution therefore establishes how strongly patient averaging changes the training objective in these allocations.

\section{Study design}
\label{sec:design}
BRACS retains its seven lesion classes and TCGA-UT its 30 eligible cancer types. Both retain the benchmark's eligibility rules, class order, and frozen 2560-dimensional Virchow2 features, formed by concatenating the CLS token and mean patch token. The three existing train--validation--test splits remain fixed. BRACS patients and TCGA-UT participants are called patients below; they are disjoint between partitions within each split \citep{queisler2026benchmarkProtocol}.

Concentrated support (\(\conc\)) and spread support (\(\spr\)) use the exact balanced training allocations evaluated by \citet{queisler2026patientShortage}: the shared concentrated allocation and the spread allocation under native class assignment. Class-specific patch counts match between support conditions within each split. The two objectives receive identical patches within each condition. Validation and test cohorts retain their natural distributions.

\begin{table}[htbp]
\centering
\renewcommand{\arraystretch}{1.15}
\begin{tabularx}{\textwidth}{@{}lXX@{}}
\toprule
Patient support & Patch-average loss & Patient-average loss \\
\midrule
Concentrated & Existing logistic baseline & Weighted fit on identical patches \\
Spread & Existing logistic baseline & Weighted fit on identical patches \\
\bottomrule
\end{tabularx}
\caption{Four cells per dataset and split compare two objectives under two support conditions, isolating the weighting intervention within each condition.}
\label{tab:design}
\end{table}

\noindent The inherited allocations contain \num{13982}--\num{27202} patches per BRACS split, with 78--89 concentrated patients and 103 spread patients. TCGA-UT contains \num{36988}--\num{52204} patches, with 671--917 concentrated patients and \num{4983} spread patients \citep{queisler2026patientShortage}. These unequal shortage doses prevent attributing differences in effect size to intrinsic dataset susceptibility. The support contrast also changes patient identities and within-patient patch contributions; it does not isolate patient count alone.

\section{Patient-average training}
\label{sec:objective}
For one training allocation, \(N\) denotes its total patch count, \(n_c\) the count for class \(c\), and \(G_c\) the number of patients contributing to that class. Patient \(i\) contributes \(m_{ic}>0\) such patches, whose indices form \(\mathcal J_{ic}\). Thus \(n_c=\sum_{i=1}^{G_c}m_{ic}\). Logistic regression predicts class probabilities from \(Wz_j+b\), where \(z_j\) is a frozen feature vector, \(W\) the coefficient matrix, and \(b\) an unpenalized intercept. Its cross-entropy for patch \(j\) is \(\ell_j=-\log p(y_j\mid z_j)\).

Ordinary training gives each patch equal weight. Patient-average training instead uses
\begin{equation}
 w_j=\frac{n_c}{G_c m_{ic}}\quad(j\in\mathcal J_{ic}),
 \qquad
 L_{\mathrm{patient}}=\frac1N\sum_{j=1}^{N}w_j\ell_j
 +\frac{\lambda}{2}\lVert W\rVert_F^2.
 \label{eq:objective}
\end{equation}
\noindent Each patient then contributes total weight \(n_c/G_c\) within class \(c\). Equivalently, the data term averages patch losses within a patient, averages patients within a class, and combines classes with their unchanged shares \(n_c/N\):
\begin{equation}
 \frac1N\sum_jw_j\ell_j
 =\sum_{c=1}^{K}\frac{n_c}{N}\frac1{G_c}
 \sum_{i=1}^{G_c}\frac1{m_{ic}}\sum_{j\in\mathcal J_{ic}}\ell_j.
 \label{eq:patient-average}
\end{equation}
\noindent Here \(K\) is the number of classes. Preserving \(n_c/N\) also preserves small integer-rounding differences in the balanced allocations. A patient contributing several classes enters each class separately; the objective does not force equal total influence across such patients irrespective of class.

The weights satisfy \(\sum_{j:y_j=c}w_j=n_c\) and \(\sum_jw_j=N\). Their mean is one, so weighting does not change the loss scale relative to regularization. If all patients contribute equal numbers of patches within each class, every weight equals one and the two objectives coincide.

\subsection{Fixed fitting procedure}
Both objectives use the deterministic multinomial logistic regression procedure of \citet{queisler2026patientShortage}, with the same feature preprocessing and numerical convergence criteria. Regularization remains fixed at the validation-selected values in Table~\ref{tab:regularization}. The analysis includes only converged fits, so the comparison concerns the weighting intervention under matched fitting conditions.

\begin{table}[htbp]
\centering
\begin{tabular}{@{}lrr@{}}
\toprule
Dataset & Concentrated \(\lambda\) & Spread \(\lambda\) \\
\midrule
BRACS & 1 & 0.1 \\
TCGA-UT & 0.01 & 0.01 \\
\bottomrule
\end{tabular}
\caption{Validation-selected regularization from \citet{queisler2026patientShortage}, held fixed for both objectives across all three splits.}
\label{tab:regularization}
\end{table}

\noindent Prediction selects the largest logit, resolving ties by the fixed class order. Validation and test patients never enter fitting or weight calculation. Freezing regularization isolates the weighting change within each support condition. On BRACS, the coverage contrast remains a comparison of procedures with different inherited regularization; it is not a fixed-regularization patient-count effect.

\subsection{Patient contributions}
The contribution distribution is characterized before fitting by \(n_c\), \(G_c\), the minimum and maximum \(m_{ic}\), and the largest patient share \(\max_i(m_{ic}/n_c)\), for every class, split, and support condition. Departure from equal influence is measured by
\begin{equation}
 D_c=\frac12\sum_{i=1}^{G_c}\left|\frac{m_{ic}}{n_c}-\frac1{G_c}\right|.
 \label{eq:influence}
\end{equation}
\noindent This is the fraction of within-class loss mass redistributed by patient averaging; zero means the objectives already coincide for that class. The measure describes inequality in patients' contributions to the training objective. Allocation and weight checks establish identical patches across objectives, matching class patch counts across support conditions, positive weights, preserved class loss mass, total weight \(N\), and unit weights under equal contributions. The benchmark's 10\% patient and 5\% slide caps constrain patch allocation; patient-average loss weights follow Equation~\eqref{eq:objective}.

\section{Evaluation}
\label{sec:evaluation}
The primary endpoint is \emph{patient-macro balanced accuracy}, with equal weighting of patients within each class as in \citet{queisler2026patientShortage}. For split \(r\), \(\mathcal P_{rc}\) contains test patients with at least one patch of true class \(c\). Their within-patient class recalls are \(q_{irc}\). The split score is
\begin{equation}
 A_r=\frac1K\sum_{c=1}^{K}\frac1{|\mathcal P_{rc}|}
 \sum_{i\in\mathcal P_{rc}}q_{irc},
 \qquad
 q_{irc}=\frac1{|\mathcal J^{\mathrm{test}}_{irc}|}
 \sum_{j\in\mathcal J^{\mathrm{test}}_{irc}}\mathbf1\{\widehat y_j=c\}.
 \label{eq:accuracy}
\end{equation}
\noindent Here \(\mathcal J^{\mathrm{test}}_{irc}\) contains that patient's evaluated patches of class \(c\). The final score averages the three split scores equally, separately for each dataset. Secondary endpoints are class-specific patient-macro recalls and patch-micro balanced accuracy, which pools patches within class before averaging classes. Probability quality and calibration are outside this experiment's scope.

The pooled score for objective \(o\in\{\mathrm{patch},\mathrm{patient}\}\) under support \(s\in\{\conc,\spr\}\) is denoted by \(A_{o,s}\). The contrasts are
\begin{align}
 W_s&=A_{\mathrm{patient},s}-A_{\mathrm{patch},s}, &&\text{weighting gain},\label{eq:gain}\\
 B_o&=A_{o,\spr}-A_{o,\conc}, &&\text{coverage benefit},\label{eq:coverage}\\
 I&=W_{\conc}-W_{\spr}=B_{\mathrm{patch}}-B_{\mathrm{patient}},
 &&\text{interaction}.\label{eq:interaction}
\end{align}
\noindent The primary questions concern \(W_{\conc}\) and \(B_{\mathrm{patient}}\). All four accuracies are reported as percentages; both weighting gains, both coverage benefits, and the interaction are expressed in percentage points. A positive interaction means the coverage gap shrinks, but could arise because spread performance deteriorates. Its interpretation therefore depends on the absolute scores and both gains.

Uncertainty is estimated with 2,000 paired Bayesian-bootstrap replicates over distinct test patients, following \citet{queisler2026patientShortage}. Each replicate assigns random positive patient weights from a \(\mathrm{Dirichlet}(1,\ldots,1)\) distribution. Each patient's weight is shared across objectives, support conditions, classes, and repeated test appearances across splits. Within-class recalls are recalculated as weighted patient means before averaging splits and forming contrasts. The 2.5th and 97.5th percentiles define individual exploratory 95\% intervals, accompanied by split-level estimates. The intervals describe patient variation conditional on the deterministic fits; they are not adjusted for multiple comparisons.

\section{Interpretation criteria}
\label{sec:interpretation}
One percentage point is the prespecified practical-relevance threshold, following \citet{queisler2026patientShortage}. A point estimate above one point is practically promising; a lower interval bound above one point supports a positive effect exceeding that threshold. Intervals crossing zero remain inconclusive about direction. An interval excluding zero but overlapping one point supports direction without establishing a practically meaningful gain. A negligible residual coverage difference requires the entire interval for \(B_{\mathrm{patient}}\) to lie within \([-1,1]\) points.

\begin{table}[htbp]
\renewcommand{\arraystretch}{1.15}
\begin{tabularx}{\textwidth}{@{}>{\raggedright\arraybackslash}p{0.34\textwidth}>{\raggedright\arraybackslash}X@{}}
\toprule
Observed pattern & Supported interpretation \\
\midrule
Positive \(W_{\conc}\), larger than \(W_{\spr}\) & Equalizing influence helps under shortage, with an additional shortage-specific benefit if the interaction is supported. \\
Positive \(W_{\conc}\) and residual \(B_{\mathrm{patient}}\) & Unequal influence contributes, but the correction leaves a coverage benefit. \\
Residual \(B_{\mathrm{patient}}\), little supported \(W_{\conc}\) & The tested correction does not explain away the coverage deficit; a null gain alone does not exclude an influence effect. \\
Positive \(W_{\conc}\), residual interval within \([-1,1]\) & The coverage gap is practically negligible under patient averaging. Spread accuracy indicates whether gap closure accompanies deterioration. \\
\bottomrule
\end{tabularx}
\caption{Prospective interpretations. Effect magnitude and uncertainty qualify every pattern; these are not experimental findings.}
\label{tab:interpretation}
\end{table}

\noindent Weighting gains are interpreted together with the patient contribution distribution. Near-zero \(D_c\) means the intervention changes little in that class, so a null result provides little evidence about strongly unequal patient influence elsewhere. Similar gains in both support conditions indicate a general weighting benefit; larger gains under concentrated support suggest a contribution specific to patient shortage. Broad or conflicting intervals leave the mechanism unresolved.

\section{Limitations}
\label{sec:limitations}
This experiment is conditional on the exact allocations and cohorts that motivated the hypothesis. Its intervals describe held-out patient variation, not variation from new training-patient draws, hyperparameter selection, or another population. Shared bootstrap weights account for repeated test patients but do not remove dependence from overlapping training and validation partitions. Fixed regularization tests this intervention at the inherited settings, not the best attainable patient-average procedure.

A residual coverage benefit cannot distinguish missing morphology, acquisition variation, patient composition, or within-patient sampling. Patient averaging changes influence without removing patch redundancy. Neither a residual deficit nor a null gain establishes an information ceiling, intrinsic inseparability, or classifier impossibility. Feature diversity, patch depth, fresh allocations, additional readouts, and other representations require separate experiments.

\bibliographystyle{plainnat}
\bibliography{references}
\end{document}
